\subsection{Two columns on the regular interval}

On a compact regular interval, the two center generators give two
positive action columns with a strictly monotone quotient. This controls
both the number of critical zeros and their crossing direction. We first
state the elementary perturbation argument used to pass from the action
to the actual return.

\begin{lemma}[Two columns and their crossing direction]
\label{rw:atlas:two-columns}
On one interval let $A,B$ have uniformly bounded $C^1$ norms, with
$|B|>b_*>0$ and $(A/B)'$ of constant sign and magnitude greater
than $r_*>0$. Put $Q=aA+bB+e$, where $a,b$ are constant,
$I=|a|+|b|>0$ and $\|e\|_{C^1}\le\eta I$.
For a sufficiently small common $\eta$, every zero is simple,
$|a|\ge cI$ there, and its derivative has sign
$aB(A/B)'$. Thus $Q$ has at most one zero, including multiplicity.
If moreover $A,B,A',B'$ and $(A/B)'$ have positive margins, then
$Q'$ has sign $a$ wherever $Q=0$ or $Q$ has sign $a$.
When $a=0$, both $Q$ and $Q'$ have sign $b$.
\end{lemma}
\begin{proof}
At a zero, $a(A/B)+b=-e/B$, so the bounded ratio and nonvanishing
$B$ give $|a|\ge cI$. Differentiating at that zero gives
\[
 Q'=aB(A/B)'+e'-B'e/B.
\]
The last two terms are $O(\eta I)$, proving the strict direction.
Two crossings with this same direction are impossible.
For the additional assertion, if $a,b$ have the same sign, use the
positive column derivatives. If their signs are opposite, at a point
where $Q=0$ or $Q$ has sign $a$,
$\operatorname{sign}(a)[a(A/B)+b]\ge-|e|/B$, again forcing
$|a|\ge cI$. In
$Q'=B'[a(A/B)+b]+aB(A/B)'+e'$ the second term has sign $a$ and dominates the possible
$O(\eta I)$ contribution of the opposite sign. The case $a=0$ follows from
the margins of $B,B'$.
\end{proof}

We keep the original central height as output. Write
\[
 \tau=\bar a+\bar c,\qquad
 \varsigma=\bar B+\bar m,\qquad \sigma=\bar c\varsigma.
\]
Intersect the upper return range $0<p<p_c$ with $F_+>0$ on
$[0,p]$ and $F_->0$ on $[0,\mathcal U_0(p)]$. Since
$\mathcal U_0$ increases, this regular limiting domain is one interval
with left endpoint zero. On it set
$\mathcal A=A_-(\mathcal U_0(p))-A_+(p)$.
A polynomial quasi-invariant of the affine upper field expresses
$\mathcal A$ in the two center generators. Its columns are positive
integrals, and two weighted-average arguments prove that their quotient
increases. The same calculation will determine the direction in which
the actual critical function crosses zero.

\begin{lemma}[The regular action columns]\label{rw:regular-columns}
For the upper orbit of flight time $T$ in
Lemma~\ref{rw:finite:upper-return}, define
\[
 J_1(T)=\int_0^T e^{-\bar ct}(1-\bar mY(t))\,\dd t,
 \qquad J_2(T)=\int_0^T e^{-\bar ct}Y(t)\,\dd t.
\]
Then $J_1=p+e^{-\bar cT}q>0$, $J_2>0$, and
\begin{equation}\label{rw:finite:quasi-invariant}
 F_-(q)-e^{\bar cT}F_+(p)
       =-e^{\bar cT}(\tau J_1+\sigma J_2).
\end{equation}
On the regular interval,
\begin{equation}\label{rw:finite:action-columns}
 \mathcal A=\tau A+\sigma C,\qquad
 A=\int_0^T\omega J_1\,\dd t,
 \quad C=\int_0^T\omega J_2\,\dd t,
 \quad \omega=\frac{p p_Te^{\bar cT}}{F_+(p)F_-(q)}>0.
\end{equation}
Both columns and their $p$ derivatives are positive, and
$(C/A)'>0$. Every nonzero constant combination of $A,C$ has at
most one zero, which is simple. At the small endpoint, uniformly
on bounded normalized coefficient sets,
\begin{equation}\label{rw:finite:small-columns}
 A=\tfrac23p^3(1+O(p)),\qquad
 C=\tfrac2{15}p^5(1+O(p)),\qquad C/A=\tfrac15p^2(1+O(p)).
\end{equation}
\end{lemma}
\begin{proof}
Introduce the polynomial
\[
 V=1+\bar aX+\bar BX^2+(2\bar B+\bar c\bar a)Y
                  +\bar c\bar BXY-\bar B\bar mY^2.
\]
Substitution of \eqref{rw:finite:principal-upper}, grouping its
constant and $Y$ terms, gives
$V'-\bar cV=-\tau(1-\bar mY)-\sigma Y$.
Its endpoint values are $F_+(p)$ and $F_-(q)$, so integration
proves \eqref{rw:finite:quasi-invariant}.
The upper arc has $Y>0$ in its interior, giving $J_2>0$.
Using $1-\bar mY=-X'$ and
$e^{-\bar ct}X=(e^{-\bar ct}Y)'$, integration by parts gives
$J_1=p+e^{-\bar cT}q$.

For monotonicity, put $u_1(T)=\int_0^T e^{-\bar ct}\,\dd t$.
Then $J_1=u_1-\bar mJ_2$.
From $Y=p(T)S-I$ and the formulas for $W,p_T,q$ above,
\[
 J_2'=p_T\int_0^T e^{-\bar ct}S(t)\,\dd t
       =e^{-\bar cT}W^2/S^2=e^{-\bar cT}q^2.
\]
Thus $J_2'/u_1'=q^2$ is strictly increasing. Since both integrals
vanish at zero, $J_2/u_1$ is its weighted average and is strictly
smaller than its terminal value. Hence
\[
 (J_2/u_1)'=\frac{u_1'}{u_1}(q^2-J_2/u_1)>0,
 \qquad
 (J_2/J_1)'=
 \frac{(J_2/u_1)'}{(1-\bar mJ_2/u_1)^2}>0.
\]
Here $1-\bar mJ_2/u_1=J_1/u_1>0$.

Differentiate $A_-(q)-A_+(p)$ and use
$q_T=(p/q)e^{\bar cT}p_T$. The result is
\[
 \mathcal A_T=
 \frac{p p_T}{F_+(p)}
       \left(\frac{e^{\bar cT}F_+(p)}{F_-(q)}-1\right)
             =\omega(\tau J_1+\sigma J_2).
\]
Both actions vanish at zero, so this proves
\eqref{rw:finite:action-columns}. The ratio $C/A$ is the weighted
average of the strictly increasing $J_2/J_1$ with positive weight
$\omega J_1$. The same averaging argument gives $(C/A)_T>0$,
hence also $(C/A)_p>0$. It proves the zero assertion directly;
positivity of the derivatives follows from the positive integrands.
Finally scale $X=p\xi$, $Y=p^2\eta$, $t=p\theta$ in the
affine field. It tends uniformly to
$\xi'=-1$, $\eta'=\xi$, whose arc is
$\eta=(1-\xi^2)/2$. Thus
$T=2p(1+O(p))$, $J_1=2p(1+O(p))$,
$J_2=\frac23p^3(1+O(p))$ and $\omega=p/2(1+O(p))$.
Integration with $\dd T=2(1+O(p))\dd p$ gives
\eqref{rw:finite:small-columns}.
\end{proof}

\begin{proposition}[The regular critical function]
\label{rw:finite:regular-count}
Fix compact regular upper and lower data, with the input bounded
away from zero. On common slightly larger tubes the closing function
is exactly
\[
 F(p)=U(p)-V(p)-\rho\nu H,
 \qquad V=V_0\circ\mathcal U_{\epsilon,\nu},
 \qquad K=\epsilon\log(V'/U').
\]
Here $U',V'<0$, and $H>h_*>0$ is input independent.
There are uniform choices of tubes, $\theta_0>0$ and a scale cutoff
such that, on any one interval contained in the tubes, if
\begin{equation}\label{rw:finite:regular-forcing}
 |\nu|\le\theta_0(|\tau|+|\sigma|),
\end{equation}
then either the field is a center or $K$ has at most one zero,
which is simple. Consequently $F$ has at most two isolated zeros,
counting multiplicity, on the positive central passing interval.
The coefficient condition holds with a factor tending to zero in
place of $\theta_0$ whenever the fibre contains a closing zero
in one of these compact regular tubes.
\end{proposition}
\begin{proof}
Use the references of Lemma~\ref{rw:finite:regular-lower}, fixed
before the input is varied. Put $q=\mathcal U_{\epsilon,\nu}(p)$
and
\[
 r_U=A_+'-\epsilon(\log\alpha_U)',\qquad
 r_V=A_-'-\epsilon(\log\alpha_V)'.
\]
They have positive margins. Direct differentiation gives the exact
identity
\begin{equation}\label{rw:finite:regular-logarithm}
 -K=A_-(q)-A_+(p)
       +\epsilon\log\frac{\alpha_U(p)r_U(p)}
        {\alpha_V(q)r_V(q)\mathcal U_{\epsilon,\nu}'(p)}.
\end{equation}
The argument of the logarithm is positive and its logarithm has
bounded selected finite jets. Thus the right side converges, with
those jets, to $\mathcal A$ at $\nu=0$, with error
$O(\epsilon(1+|\log\epsilon|))$; changing $\nu$ adds $O(\nu)$.

The real period domains of Section~5 contain each compact regular
center tube under consideration. On the second component, the
separated coordinate at $Y=0$ is $\chi=-p$ and
$A_C(-p)=F_+(p)+O(\epsilon)$ throughout the bounded interval.
Its positive margin places the entire arc in the actual center
domain. On the reversible component the compact upper arc lies
strictly inside the central potential well. Hence in either case
the complete central maps agree. Since the reference values vanish
at $\nu=0$, $K=0$ there. The zero-scale action has the same two
identities by \eqref{rw:finite:action-columns}.

Choose a common larger regular coefficient domain near each
center tube and apply Proposition~\ref{rw:generator-division} to
the difference between \eqref{rw:finite:regular-logarithm} and
$\mathcal A$. Two additional coefficient derivatives suffice.
Away from the center set the generator norm has a positive margin,
so the error can instead be divided by that norm. On the resulting
compact cover we have
\begin{equation}\label{rw:finite:perturbed-columns}
 -K=\tau A_\epsilon+\sigma C_\epsilon+\nu E_\epsilon,
 \qquad (A_\epsilon,C_\epsilon)\longrightarrow(A,C)
                  \quad\hbox{in }C^2_p,
 \quad \|E_\epsilon\|_{C^2_p}\le C.
\end{equation}
The common center domain is essential here: it allows division of
the relative logarithmic error by the exact generators, with the
stated input-derivative bounds.

On the larger tubes both actions have a common lower bound $S_*>0$.
The amplitude and upper-map estimates bound the required absolute
coefficient derivatives by $C\epsilon^{-M}e^{-S_*/\epsilon}$,
also on the small $\nu$ box after reducing $S_*$. The same center
division preserves this weight. Lemma~\ref{rw:centers:relative-forcing},
with $I=|\tau|+|\sigma|$, therefore gives
$|\nu|=o(1)I$ at a closing zero; for $I=0$ it gives the center
identity and no isolated zero.

Choose $\theta_0$ and then the scale so that
\eqref{rw:finite:perturbed-columns} satisfies
Lemma~\ref{rw:atlas:two-columns}, with columns $(C,A)$ and
coefficients $(\sigma,\tau)$. The strict quotient margin of
Lemma~\ref{rw:regular-columns} gives at most one simple zero of $K$.
Lemma~\ref{rw:zero-calculus} gives the stated count for $F$.
Since $L_+(p)$ and $L_-(q(p))$ decrease, their positive sets are
initial subintervals. Intersecting them with the regular tube gives
one positive-central passing interval, to which the count applies.
\end{proof}

\begin{lemma}[Direction on the regular part]
\label{rw:finite:regular-direction}
Under the preceding choices, if $\sigma\ne0$, then at every point
where $K=0$ or $K$ has sign $-\sigma$, its derivative in increasing
$p$ has sign $-\sigma$. After its possible zero it therefore keeps
this sign and increases in absolute value on the rest of the regular
passing interval. If $\sigma=0$, $\tau\ne0$, both $K$ and $K'$
have sign $-\tau$.
\end{lemma}
\begin{proof}
The positive margins of $A,A',C,C'$ and $(C/A)'$ on each compact
regular tube allow the last clause of Lemma~\ref{rw:atlas:two-columns}
to be applied to $-K$ in \eqref{rw:finite:perturbed-columns}.
At a zero, $-K$ crosses into the side of sign $\sigma$, and its
derivative keeps that sign there. Thus it stays on that side. A regular
input change preserves this direction for the same original-central
function $K$.
\end{proof}
